\subsubsection{Digital Images as Arrays}\label{sec:digital-images-as-arrays}

Until now, we have discussed what Computer Vision does. The next question is: \emph{\textbf{how does a computer actually represent an image?}} Humans perceive an image as a continuous visual scene, but computers cannot directly understand colors, shapes, or objects. Before any neural network can process an image, the image must first be represented numerically. This is the purpose of this section.

\highspace
\begin{flushleft}
    \textcolor{Green3}{\faIcon{question-circle} \textbf{Why represent images as arrays?}}
\end{flushleft}
A neural network only knows how to perform mathematical operations. It cannot process concepts like ``\emph{red car}'' or ``\emph{person sitting}''. Instead, it receives a large collection of numbers. Therefore, every digital image must be converted into a \textbf{multidimensional array}. The network will later learn how to associate patterns in these numbers with semantic concepts.

\highspace
In general, an image is composed of a finite number of tiny elements called \textbf{pixels} (picture elements). Each pixel stores information about the light intensity reaching that location. If the image is grayscale, every pixel stores only \textbf{one number}. If the image is colored, every pixel stores \textbf{three numbers}.
\begin{itemize}
    \item \important{Grayscale images}. Suppose we have a grayscale image. It can be represented as:
    \begin{equation*}
        I \in \mathbb{R}^{R \times C}
    \end{equation*}
    Where $R$ is the number of rows (height) and $C$ is the number of columns (width). Each element $I\left(r,c\right)$ contains the brightness of one pixel.

    For example, consider the following $3 \times 3$ grayscale image:
    \begin{equation*}
        I = \begin{bmatrix}
            10 & 35 & 80 \\
            50 & 120 & 200 \\
            15 & 90 & 255
        \end{bmatrix}
    \end{equation*}
    Means that dark pixels correspond to small numbers, while bright pixels correspond to large numbers. So a grayscale image is simply a \textbf{2-dimensional matrix}.


    \item \important{RGB images}. Most images are not grayscale. Instead, they are represented using the \textbf{RGB color model}. RGB stands for \emph{Red}, \emph{Green}, and \emph{Blue}. Every color can be generated by combining different amounts of these three primary colors. Each channel measures \textbf{how much of that color is present} at every pixel. When combined, they reconstruct the original color image.
    
    Mathematically, instead of one matrix, an RGB image consists of \textbf{three matrices}: $R$, $G$, and $B$. Each having size $R \times C$. Therefore, the complete image becomes a \textbf{3-dimensional array}:
    \begin{equation*}
        I \in \mathbb{R}^{R \times C \times 3}
    \end{equation*}
    The third dimension stores the three color channels.
\end{itemize}

\highspace
\begin{flushleft}
    \textcolor{Green3}{\faIcon{tools} \textbf{Pixel values}}
\end{flushleft}
Images are usually stored using \textbf{8 bits per color channel}, so each channel value is an integer in the range \textbf{0-255}. Since $2^{8} = 256$, each channel has 256 possible values: $0$ represents no intensity (black), while $255$ represents full intensity (white). Therefore, a pixel in an RGB image can be represented as a triplet of integers:
\begin{equation*}
    \text{pixel} = \left(R, G, B\right) \in \left[0, 255\right]^{3}
\end{equation*}
For example, the pixel $\left(253, 5, 6\right)$ means that the red channel is almost at full intensity, while the green and blue channels are very low. This pixel will appear as a bright red color.

\highspace
\textcolor{Green3}{\faIcon{question-circle} \textbf{If pixel values are integers between $0$ and $255$, why are we saying the image belongs to $\mathbb{R}$ (real numbers)?}} When the image is read from disk, the pixel values are indeed integers (\texttt{uint8}). However, before entering a neural network, these integers are almost always converted to floating-point numbers (\texttt{float32}), often normalized to the range $[0, 1]$ or $[-1, 1]$. This is because later computations (matrix multiplications, convolutions, gradients, etc.) are all performed on \textbf{real-valued tensors}.

\highspace
\begin{flushleft}
    \textcolor{DarkOrange3}{\faIcon{balance-scale} \textbf{Images in memory vs. on disk}}
\end{flushleft}
When loaded into memory, \hl{images occupy much more space than they do on disk because image files (e.g., JPEG) are typically compressed.} For example, a JPEG image might occupy only a few hundred kilobytes on disk. After decompression, it becomes a full array containing every pixel, which is much larger. This distinction is important because \hl{neural networks always operate on \textbf{decompressed numerical array}}, not on the compressed JPEG representation.

\highspace
In general, a neural network never receives an image as a ``picture''. Instead, it receives $I \in \mathbb{R}^{R \times C \times 3}$, which is simply a \textbf{3-dimensional tensor of numbers}. Everything the network learns (e.g., edges, textures, objects, faces, animals, etc.) is ultimately learned by discovering patterns within these numerical arrays.